
qu'un
nombre fini de points singuliers, r\'eduction justifi\'ee par la formule de projection X 4.4.4 et par 4.5 (ii). Il existe donc une courbe ferm\'ee int\`egre @@M0@@ sur @@M1@@ telle que @@M2@@, et @@M3@@ est localement principale en dehors d'un ensemble fini @@M4@@. Soit @@M5@@ l'\'eclatement de l'id\'eal @@M6@@ de @@M7@@. L'id\'eal @@M8@@ est inversible, et il d\'efinit un sous-sch\'ema de @@M9@@, qui est la r\'eunion de deux sous-sch\'emas @@M10@@, @@M11@@ tels que @@M12@@ et @@M13@@. Gr\^ace \`a la formule de projection X 4.4.4, on a que @@M14@@ est \'egal \`a @@M15@@, et on peut remplacer @@M16@@ par @@M17@@ ; on est donc ramen\'e au cas pr\'ec\'edent, et la d\'emonstration de 7.1 est achev\'ee.

\textbf{Remarque 7.2.} On ignore si une surface irr\'eductible, propre (sur un corps alg\'ebriquement clos), qui poss\`ede un faisceau inversible @@M18@@ tel que @@M19@@, est n\'ecessairement projective. On ignore \'egalement si pour toute surface irr\'eductible, propre, la forme d'intersection sur @@M20@@ est non-d\'eg\'en\'er\'ee ; il r\'esulte de 7.1 que lorsqu'elle est d\'eg\'en\'er\'ee, elle est n\'ecessairement n\'egative (non-d\'efinie).

\textbf{Corollaire 7.3.} \emph{Sous les conditions de 7.1, soit @@M21@@ un faisceau inversible sur @@M22@@ qui n'est pas num\'eriquement proportionnel \`a @@M23@@, et soient @@M24@@ et @@M25@@. Alors, @@M26@@.}

En effet, soit @@M27@@. Comme @@M28@@, d'apr\`es 7.1, @@M29@@ ; d'o\`u l'assertion.

